\section{Four-input comparisons and the pentagon}
\label{sec:four-leaf-pentagon}

\subsection{Six divisors and five bracketings}
\label{fp:sec:four-divisors}

Four positions have six pairwise collision divisors.  A binary tree
selects three of them when each merger retains its rightmost position.
The five selections admit a single coefficient diagram.  Its maps
satisfy a pentagon, but its projections are equivalences only on
specified open subsets.  The distinction is detected by a single pole
class at the generic point of each omitted divisor.

In this section \(k\) is any field.  Put
\[
 a=z_1-z_2,\quad b=z_2-z_3,\quad c=z_3-z_4,\qquad R=k[a,b,c],
\]
\[
 d=a+b=z_1-z_3,\qquad e=b+c=z_2-z_4,\qquad
 f=a+b+c=z_1-z_4.
\]
Write \(\mathcal D=\{a,b,c,d,e,f\}\).
The symbol \(d\) denotes the spatial function in this display.
Write \(\partial\) for the differential of a complex.
The six functions are pairwise nonassociate irreducible polynomials,
including in characteristics two and three.  Fix the factor order
\(a,b,c,d,e,f\).  All permutations of tensor factors below use the
Koszul sign.  State coefficients are placed after interval coefficients.

For \(\delta\in\mathcal D\), use the interval algebra
\(A_\delta=\mathcal I_R(\delta)\) of
\eqref{sb:eq:general-interval}, with augmentation \(p_\delta\)
and unit map \(i_\delta:R\to A_\delta\).
The differential is denoted by \(\partial\).

Let \(D_R\) be a unital dg \(R\)-algebra.  Its multiplication
may be noncommutative.  For \(J\subseteq\mathcal D\), define
\begin{equation}\label{fp:eq:four-common-carrier}
 R_J=R[\delta^{-1}:\delta\in J],\qquad
 \mathcal C_J(D_R)=R_J\otimes_R
       \left(\bigotimes_{\delta\in\mathcal D}A_\delta\right)
       \otimes_R D_R.
\end{equation}
The inclusions \(J\subseteq J'\) give scalar-extension maps.
They compose strictly.  Thus \(\mathcal C_\bullet(D_R)\)
is a diagram indexed by the subsets of the six divisors.
Its initial object keeps the endpoint ring \(R\).
In coefficient degree \(r\), its summands have overlap ring
\[
 R_J[\delta^{-1}:\delta\in I],\qquad
 I\subseteq\mathcal D,\quad |I|=r,
\]
with one endpoint choice for each \(\delta\notin I\).
This specifies every module, including simultaneous overlaps.
No completion or intersection of Laurent-series fields is used.

Number the bracketings and their selected sets as follows:
\begin{equation}\label{fp:eq:four-tree-sets}
\begin{array}{c|c|c}
 T&\text{bracketing}&B_T\\\hline
 0&(((1,2),3),4)&\{a,b,c\}\\
 1&((1,(2,3)),4)&\{b,c,d\}\\
 2&(1,((2,3),4))&\{b,c,f\}\\
 3&(1,(2,(3,4)))&\{c,e,f\}\\
 4&((1,2),(3,4))&\{a,c,e\}
\end{array}
\end{equation}
Set \(E_T(D_R)=\bigotimes_{\delta\in B_T}A_\delta
\otimes_R D_R\), in the fixed order.
There are strict dg algebra maps
\[
 j_T:E_T(D_R)\longrightarrow\mathcal C_\varnothing(D_R),
 \qquad
 q_T:\mathcal C_\varnothing(D_R)\longrightarrow E_T(D_R).
\]
The map \(j_T\) inserts units in missing factors.  The map
\(q_T\) evaluates those factors at their first endpoints.
They preserve the state factor and satisfy \(q_Tj_T=1\).
The same formulas commute with every scalar extension in
\eqref{fp:eq:four-common-carrier}.

\begin{proposition}[The retained divisors and the tree projections]
\label{fp:prop:four-full-carrier}
For every \(\delta\in\mathcal D\), the initial scalar carrier
\(\mathcal C_\varnothing(R)\) has a nonzero degree-one class
at the prime \((\delta)\).  That class is represented by
\(\theta_\delta/\delta\), with units in the other factors.
For \(J_T=\mathcal D\setminus B_T\), the maps \(j_T,q_T\)
become mutually inverse chain homotopy equivalences over \(R_{J_T}\),
for arbitrary \(D_R\).  They need not be equivalences over \(R\).
\end{proposition}

\begin{proof}
At \((\delta)\), every other collision function is a unit.
For an invertible \(\gamma\), set
\[
 h_\gamma(\ell\theta_\gamma)=\ell e_1^\gamma,
 \qquad h_\gamma(A_\gamma^0)=0,\qquad
 F_\gamma=i_\gamma p_\gamma.
\]
A direct calculation gives
\(\partial h_\gamma+h_\gamma\partial=1-F_\gamma\),
\(h_\gamma^2=0\), and
\(p_\gamma h_\gamma=h_\gamma i_\gamma=0\).
Tensor these maps with the other factors, with the tensor signs.
Contracting the five invertible intervals identifies the local
cohomology with that of \(A_\delta\).  In degree one it is
\[
 R_{(\delta)}[\delta^{-1}]/R_{(\delta)}.
\]
The indicated pole class is nonzero.

More generally, to remove invertible intervals
\(\gamma_1,\ldots,\gamma_m\) in a specified order, use
\begin{equation}\label{fp:eq:four-tensor-contraction}
 h=h_{\gamma_1}+F_{\gamma_1}h_{\gamma_2}
       +\cdots+F_{\gamma_1}\cdots F_{\gamma_{m-1}}h_{\gamma_m}.
\end{equation}
Each \(h_\gamma\) includes its tensor sign and acts trivially
on \(D_R\).  Differentiation telescopes to
\(1-F_{\gamma_1}\cdots F_{\gamma_m}\).
For \(\gamma_i\in J_T\), the last product is \(j_Tq_T\).
Together with \(q_Tj_T=1\), this proves the second assertion.
The construction is an explicit chain contraction and needs no
flatness assumption on \(D_R\).
\end{proof}

The initial carrier retains all six generic collisions.  The whole
diagram records exactly where a projection discards their pole classes.
Replacing the diagram by its last object would discard this information.
The proposition asserts retention for the scalar coefficient carrier.
For an arbitrary state algebra, tensoring can annihilate classes.

\subsection{The maximal common open and its pentagon}
\label{fp:sec:four-coefficient-pentagon}

The adjacent rotations are
\[
 0\longrightarrow1\longrightarrow2\longrightarrow3,
 \qquad 0\longrightarrow4\longrightarrow3.
\]
For any arrow \(T\to U\) in these two paths define
\[
 r_{UT}=q_Uj_T:E_T(D_R)\longrightarrow E_U(D_R).
\]
The displayed maps are defined over \(R\).  On adjacent trees
they preserve the two common intervals, remove one interval by
first-endpoint evaluation, and insert one unit:
\begin{equation}\label{fp:eq:four-adjacent-maps}
\begin{array}{c|cc|c}
 r_{UT}&\text{removed}&\text{inserted}&\text{preserved}\\\hline
 r_{10}&a&d&b,c\\
 r_{21}&d&f&b,c\\
 r_{32}&b&e&c,f\\
 r_{40}&b&e&a,c\\
 r_{34}&a&f&c,e
\end{array}
\end{equation}
They are the endpoint rotations of the actual subtree coordinates.
For example, \(1\to2\) uses the three block representatives
\(z_1,z_3,z_4\), with differences \(d,c,f=d+c\).
The map \(4\to3\) uses \(z_1,z_2,z_4\), with differences
\(a,e,f=a+e\).

\begin{theorem}[The coefficient pentagon and its exact domain]
\label{fp:thm:four-coefficient-pentagon}
For every dg \(R\)-algebra \(D_R\), there is an equality
of strict algebra chain maps over \(R\):
\begin{equation}\label{fp:eq:four-pentagon-coefficients}
 r_{32}r_{21}r_{10}=r_{34}r_{40}:E_0(D_R)\longrightarrow E_3(D_R).
\end{equation}
Both sides preserve \(A_c\) and \(D_R\), evaluate \(A_a,A_b\),
and insert the units of \(A_e,A_f\).

Put
\begin{equation}\label{fp:eq:four-maximal-open}
 S=R[a^{-1},b^{-1},d^{-1},e^{-1},f^{-1}].
\end{equation}
Over \(S\), each adjacent map is a chain homotopy equivalence.
Among nonzero localizations of \(R\), this is the least localization
on which all five scalar tree complexes can be quasi-isomorphic.
In particular, no such common-open equivalence retains any of the
five generic divisors \(a,b,d,e,f\).  It retains \(c=0\).
\end{theorem}

\begin{proof}
Follow each original tensor factor through either path.  The long
path removes \(a\), inserts and then removes \(d\), and removes
\(b\).  The short path removes \(b\) and then \(a\).
In both cases the only surviving original interval is \(c\), and
the final \(e,f\) factors are units.  Endpoint evaluation after
unit insertion is the identity on scalars.  The fixed factor order
and the Koszul convention give identical maps, proving
\eqref{fp:eq:four-pentagon-coefficients}.

For one adjacent pair, write its complexes as the tensor of their
common factors, the changed factor, and the state algebra.  The
removed and inserted coordinates are units over \(S\).
Their explicit homotopies in Proposition~\ref{fp:prop:four-full-carrier}
show that the reverse endpoint rotation is a homotopy inverse.

For necessity, let \(R'\) be a nonzero localization of \(R\).
If a function \(\delta\) in \(\{a,b,d,e,f\}\) is not a unit
in \(R'\), its prime ideal survives.  Indeed, a multiplicative
set meeting \((\delta)\) would invert a multiple of \(\delta\)
and hence \(\delta\) itself.  Some \(B_T\) contains \(\delta\)
and another \(B_U\) omits it, by \eqref{fp:eq:four-tree-sets}.
Localizing further at \((\delta)\), the first tree has
\(H^1=R_{(\delta)}[\delta^{-1}]/R_{(\delta)}\ne0\).
Every interval of the second tree contracts, so its \(H^1\) is zero.
Localization preserves cycles and boundaries, since denominators
can be cleared.  Thus the two scalar complexes cannot be
quasi-isomorphic over \(R'\).  All five functions must be units,
so the map \(R\to R'\) factors through \(S\).
Conversely, the preceding contractions work over \(S\).
Finally \(c^{-1}\notin S\), so the class
\(\theta_c/c\) survives in \(H^1(A_c\otimes_R S)\).
\end{proof}

For an explicit obstruction, compare \(T=0\) and \(T=1\)
at \((a)\).  The class \(\theta_a/a\) in the first tree
is nonzero.  The intervals \(b,c,d\) in the second tree
all contract.  Any equivalence preserving this generic boundary
would have to preserve a class in a zero cohomology group.
Adding state factors or further comparison maps cannot remove this
contradiction while retaining the scalar complexes as equivalent fibers.

\subsection{Four polynomial state inputs}
\label{fp:sec:four-state-pentagon}

The state operation uses all six pair distances, even when a tree
selects only three interval factors.  This fact fixes its domain.
Let \(\mathsf V\) be the one-species algebra
\eqref{fs:eq:state-algebra}, with its auxiliary weight and Hasse
translation.  Write \(P_{rs}^{ij}=P_{ij}^{rs}\) and, in this section,
\(D_{ij}=\mathsf D_{ij}\) for the rational odd operator of
\eqref{fs:eq:rational-pair-operator}.  These operators act on the
original ordered tensor product of polynomial states.  They anticommute
and square to zero, by the left odd derivative relations, also in
characteristic two.

For two inputs at separation \(v\), the binary operation is
\[
 Y_v=E_v(1+\theta_v D_{12}),\qquad
 E_v(\alpha\otimes\beta)=(\tau_v\alpha)\beta.
\]
In this formula \(D_{12}\) uses separation \(v\).
Set \(z_4=0,z_3=c,z_2=e,z_1=f\), and write
\[
 E(\alpha_1\otimes\alpha_2\otimes\alpha_3\otimes\alpha_4)
  =(\tau_f\alpha_1)(\tau_e\alpha_2)(\tau_c\alpha_3)\alpha_4.
\]
Products of coefficient operators use the graded tensor rule.
Thus moving an odd \(D_{ij}\) across an interval form changes
its sign.  These conventions retain multiple contractions.

For a state algebra \(V\) over \(k\), write
\(\mathcal C_J(V)=\mathcal C_J(R\otimes_kV)\).
For a fixed tree, extend the endpoint ring to \(R_{J_T}\),
where \(J_T=\mathcal D\setminus B_T\).  The state output is
understood weightwise in
\[
 \bigcup_{N\in\mathbb N^{\mathcal D}}
 \prod_{w\ge0}F_N\bigl(E_T(R_{J_T})\otimes_k\mathsf V_w\bigr).
\]
Here \(F_N\) uses the intersection of each coefficient module
with \(\prod_{\delta\in\mathcal D}\delta^{-N_\delta}R\).
The same pole tuple bounds every output weight.  The input states
are polynomials, not arbitrary elements of the product of weights.

\begin{proposition}[The actual four-input operations]
\label{fp:prop:four-state-formulas}
The five binary iterations have the following rational representatives:
\begin{equation}\label{fp:eq:four-state-formulas}
\begin{aligned}
 Y_0&=E(1+\theta_c(D_{14}+D_{24}+D_{34}))
         (1+\theta_b(D_{13}+D_{23}))(1+\theta_aD_{12}),\\
 Y_1&=E(1+\theta_c(D_{14}+D_{24}+D_{34}))
         (1+\theta_d(D_{12}+D_{13}))(1+\theta_bD_{23}),\\
 Y_2&=E(1+\theta_f(D_{12}+D_{13}+D_{14}))
         (1+\theta_c(D_{24}+D_{34}))(1+\theta_bD_{23}),\\
 Y_3&=E(1+\theta_f(D_{12}+D_{13}+D_{14}))
         (1+\theta_e(D_{23}+D_{24}))(1+\theta_cD_{34}),\\
 Y_4&=E(1+\theta_e(D_{13}+D_{14}+D_{23}+D_{24}))
         (1+\theta_cD_{34})(1+\theta_aD_{12}).
\end{aligned}
\end{equation}
Each is a degree-zero chain map on \(\mathsf V^{\otimes4}\),
with the stated output domain over \(R_{J_T}\).
There are at most three contractions.
For every tree, \(R_{J_T}\) is the least endpoint localization
supporting these rational representatives for all polynomial inputs.
\end{proposition}

\begin{proof}
At an internal vertex, the derivative of each merged state is the
sum of derivatives on its original input factors.  A contraction
across the two children therefore sums over pairs with one leaf
in each child.  A fixed pair occurs at its unique lowest common
ancestor.  The sets at the three vertices give precisely the three
operator factors in each line of \eqref{fp:eq:four-state-formulas}.

The separation in each coefficient must be the original leaf distance.
The identities that transport it through a merger are
\[
 \sum_{m\ge0}\binom{r+m}{r}t^mK_{r+m,s}(x)=K_{rs}(x+t),
\]
\[
 \sum_{m\ge0}\binom{s+m}{s}x^mK_{r,s+m}(u)=K_{rs}(u-x).
\]
For the first, use
\(\binom{r+m}{r}\binom{r+m+s}{r+m}
=\binom{r+s}{r}\binom{r+s+m}{m}\).
The negative-binomial expansion then gives the equality, including
its sign.  The second identity follows from the same equality
without the alternating power of the translating variable.
They hold over the integers in the respective formal domains.
Applying them at each of the three vertices and using the translation
law proves the formulas, with all double and triple terms.
For the two disjoint inner vertices of \(T=4\), the even operators
\(\theta_aD_{12}\) and \(\theta_cD_{34}\) commute.
The exchange of forms cancels the exchange of odd derivatives.

Only finitely many derivative indices act nontrivially on polynomial
inputs before translation.  Each vertex contributes at most one
contraction.  Each fixed output weight has finitely many translation
terms, and translation adds only polynomial spatial coefficients.
The denominators of a contraction are powers of its leaf distance.
A distance in \(J_T\) is already invertible at every endpoint.
A distance in \(B_T\) labels its own vertex, so its denominator
appears only together with that vertex's form.  Products belong to
the corresponding simultaneous-overlap module.  This proves the
output and pole assertions.  All terms are closed.

To prove minimality, take inputs \(x_0\) and \(\eta_0\)
at a chosen pair \(i<j\), with units at the other leaves.
The only contraction is
\(\theta_\gamma/(z_i-z_j)\), where \(\gamma\in B_T\)
is the coordinate of the pair's lowest common ancestor.
If \(z_i-z_j\notin B_T\), its inverse can belong to the
\(\gamma\)-overlap ring of an endpoint localization only if
that leaf distance is already a unit there.  Indeed, its prime
ideal survives every inversion of the distinct irreducible
\(\gamma\).  Varying the pair forces every member of \(J_T\)
to be a unit.  The preceding construction proves sufficiency.
\end{proof}

\begin{example}[Three simultaneous contractions]
For inputs \((x_0,x_0\eta_0,x_0\eta_0,\eta_0)\), the fully
contracted terms, in the fixed interval order, are
\[
 \begin{array}{c|c}
 T&\text{coefficient of the scalar state}\\\hline
 0&\theta_a\theta_b\theta_c/(abc)\\
 1&f\theta_b\theta_c\theta_d/(acde)\\
 2&f\theta_b\theta_c\theta_f/(acde)\\
 3&-\theta_c\theta_e\theta_f/(abc)\\
 4&-\theta_a\theta_c\theta_e/(abc)
 \end{array}
\]
For \(T=0\), the successive nonzero contractions are
\(P^{12}_{00},P^{23}_{00},P^{34}_{00}\), each with coefficient
\(+1\) on the remaining polynomial.  For \(T=1,2\), the
signed sum is \(1/(abc)-1/(bde)=f/(acde)\).
For \(T=3,4\), left differentiation of the two odd factors
remaining at the last inner vertex gives the minus sign.
Hasse translation acts trivially on the final scalar state.
These integer identities retain their meanings in every characteristic.
\end{example}

Here are explicit formal targets for these rational representatives.
In each row, list the root variable first, followed by internal
variables from the root toward the leaves:
\[
 \begin{array}{c|c}
 T&\text{iterated Laurent field}\\\hline
 0&k((c))((b))((a))\\
 1&k((c))((d))((b))\\
 2&k((f))((c))((b))\\
 3&k((f))((e))((c))\\
 4&k((e))((c))((a))
 \end{array}
\]
For each row the three displayed variables are algebraically independent.
Every leaf distance is a nonzero linear expression in them, so its
inverse has a unique Laurent expansion in that field.  For example,
in row 2, \(a=f-c-b\), \(d=f-c\), and \(e=b+c\).
At every ancestor, descendants are expanded as the smaller variables.
The two disjoint inner operations in row 4 admit the stated order
because they commute as proved above.  Each rational embedding is
injective.  These targets are kept distinct, and comparisons below
are made before expansion.

\subsubsection{A common state operation before localization}

The complete coefficient carrier supports the closed degree-zero map
\begin{equation}\label{fp:eq:four-universal-state}
 \mathcal U=E\prod_{1\le i<j\le4}
       (1+\theta_{z_i-z_j}D_{ij})
 :\mathsf V^{\otimes4}\longrightarrow
       \widehat{\mathcal C_\varnothing(\mathsf V)}.
\end{equation}
The hat means the uniformly pole-bounded product of output weights
specified above.  There are at most six contractions.  The factors
commute: exchanging their forms and their odd derivatives gives two
opposite signs.  Thus any fixed order of the six pairs gives the
same map.  A contraction denominator lies in its own pair-overlap
module, so the operation is defined over \(R\) with all six
divisors retained.

For each tree, let \(v_T(i,j)\in B_T\) be the coordinate
of the lowest common ancestor of leaves \(i,j\).  Over \(R_{J_T}\),
define an algebra chain map from a single pair interval to that vertex
interval by
\[
 e_\nu^{z_i-z_j}\longmapsto e_\nu^{v_T(i,j)},\qquad
 \ell\theta_{z_i-z_j}\longmapsto\ell\theta_{v_T(i,j)}.
\]
It fixes every spatial scalar.  If the pair interval is selected by
the tree, then \(z_i-z_j=v_T(i,j)\) and this is the identity.
Otherwise its distance is inverted in \(R_{J_T}\), so all its
coefficients already belong to that endpoint ring.
Regroup the six factors by vertex, using the Koszul signs, and
multiply each group in the original order of \(a,b,c,d,e,f\).
This defines a degree-zero chain map
\begin{equation}\label{fp:eq:four-grouping-map}
 \mu_T:\mathcal C_{J_T}(D_R)\longrightarrow
             E_T(D_R)\otimes_R R_{J_T}.
\end{equation}
It acts identically on \(D_R\).  Multiplication in a dg algebra
is a chain map, so every stage of this definition is a chain map.
It preserves the unit and first-endpoint evaluation, and
\(\mu_Tj_T=1\).

The map \(\mu_T\) is generally not multiplicative.  Two
source factors supercommute.  If both map to a vertex \(v\),
one may choose \(e_1\) in the first and \(\theta\) in
the second.  Their images satisfy \(e_1^v\theta_v=0\) but
\(\theta_ve_1^v=\theta_v\).  A multiplicative map could
not preserve their source commutation relation.

\begin{theorem}[The state carrier with all six divisors]
\label{fp:thm:four-common-state}
Over every field, \(\mu_T\mathcal U=Y_T\) over \(R_{J_T}\).
The map \(\mu_T\) is a chain homotopy equivalence with inverse
\(j_T\).  It preserves the full rational coefficient and its
six-component pole tuple.

Let \(\widetilde E_T(D_R)\) be the tensor of the selected
intervals, followed by all missing intervals, followed by \(D_R\).
The signed reordering
\(\Phi_T:\widetilde E_T(D_R)\to\mathcal C_\varnothing(D_R)\)
is a strict dg algebra isomorphism.  The maps
\[
 \rho_{UT}=\Phi_U^{-1}\Phi_T,\qquad
 \widetilde Y_T=\Phi_T^{-1}\mathcal U
\]
satisfy a strict pentagon over \(R\), with all six divisors retained,
and \(\rho_{UT}\widetilde Y_T=\widetilde Y_U\).
Their comparisons to the native operations are
\(\mu_T\Phi_T\) over the stated \(R_{J_T}\).
\end{theorem}

\begin{proof}
Regroup \eqref{fp:eq:four-universal-state} by lowest common ancestor.
In the image of \(\mu_T\), two contraction forms at the same
vertex multiply to zero.  The terms at that vertex therefore reduce to
\(1+\theta_v\sum_{v_T(i,j)=v}D_{ij}\).
The products over the three vertices give precisely
\eqref{fp:eq:four-state-formulas}.  Regrouping introduces the same
Koszul signs on both sides.  This proves the operation identity
on the actual polynomial state tensors.

Let \(h\) be \eqref{fp:eq:four-tensor-contraction} on the common
carrier, removing the factors in \(J_T\), and let \(q=q_T\),
\(j=j_T\), \(\mu=\mu_T\).  Its side identities are
\(h^2=hj=qh=0\).  For example, two distinct summands in
\(h^2\) contain either \(h_\gamma F_\gamma\) or
\(F_\gamma h_\gamma\), both zero.
Define
\begin{equation}\label{fp:eq:four-grouping-contraction}
 H=(1-j\mu)h.
\end{equation}
Since \(\mu j=1\), direct multiplication gives
\[
 \partial H+H\partial=(1-j\mu)(1-jq)=1-j\mu,
 \qquad H^2=Hj=\mu H=0.
\]
For the square, use \(h^2=hj=0\) in
\((h-j\mu h)^2\).  Thus \((j,\mu,H)\) is an explicit
chain contraction.  Every factor map, multiplication, and homotopy
leaves spatial scalars fixed.  Products add their pole tuples.
This proves the assertions about domains and poles.

Finally \(\Phi_U\rho_{UT}=\Phi_T\).  Along any path the
successive reorderings cancel, proving the pentagon and the equality
of state operations.  The initial coefficient carrier retains the
six generic classes by Proposition~\ref{fp:prop:four-full-carrier}.
No endpoint localization is used for these enriched presentations.
The comparison to a native tree uses exactly the localization already
required to define \(\mu_T\).
\end{proof}

The enriched presentations contain the three additional interval
factors.  Their maps to the native trees are not equivalences over
\(R\).  The generic-divisor obstruction proves this distinction.

Over \(S\), let \(p_T:E_T(S)\to A_c\otimes_R S\)
evaluate all intervals other than \(c\), and let \(i_T\)
insert their units.  Define
\[
 C_0=C_1=D_{14}+D_{24}+D_{34},\qquad
 C_2=D_{24}+D_{34},\qquad C_3=C_4=D_{34},
\]
\[
 Z=E(1+\theta_cD_{34}).
\]
Then \(p_TY_T=E(1+\theta_cC_T)\).
All coefficients of \(C_T-D_{34}\) lie in \(S\), since
the only remaining denominators are powers of \(e\) and \(f\).
Let \(h_T\) be \eqref{fp:eq:four-tensor-contraction}, removing
\(B_T\setminus\{c\}\) in the fixed order.  It acts on the
interval coefficients and leaves the state map unchanged.
Set
\begin{equation}\label{fp:eq:four-state-potentials}
 G_T=h_TY_T+i_T\bigl(e_1^cE(C_T-D_{34})\bigr),\qquad
 H_{UT}=r_{UT}G_T-G_U.
\end{equation}
These are fully specified finite operator formulas.  The map \(G_T\)
has degree minus one into the \(T\)-output.  The map \(H_{UT}\)
has degree minus one into the \(U\)-output.

\begin{theorem}[A pentagon of state homotopies]
\label{fp:thm:four-state-pentagon}
For polynomial inputs over every field, the maps in
\eqref{fp:eq:four-state-potentials} satisfy
\[
 \partial G_T=Y_T-i_TZ,\qquad
 \partial H_{UT}=r_{UT}Y_T-Y_U.
\]
Their pentagon boundary is zero as an actual degree-minus-one map:
\begin{equation}\label{fp:eq:four-state-pentagon}
 r_{32}r_{21}H_{10}+r_{32}H_{21}+H_{32}
       =r_{34}H_{40}+H_{34}.
\end{equation}
All maps retain the \(c\)-collision and preserve the total
six-component pole bound.  No scalar-state quotient is used.
\end{theorem}

\begin{proof}
Every \(Y_T\) is closed, and
\(\partial h_T+h_T\partial=1-i_Tp_T\).
Also
\(\partial(e_1^cE(C_T-D_{34}))=\theta_cE(C_T-D_{34})\).
The latter primitive is legal because its coefficients lie in \(S\).
Combining these identities gives
\[
 \partial G_T=Y_T-i_Tp_TY_T+i_T(p_TY_T-Z)=Y_T-i_TZ.
\]
Each adjacent map preserves \(c\), the unit factors, and the
state coefficients, so \(r_{UT}i_T=i_U\).
Applying \(\partial\) to \(H_{UT}\) proves its chain equation.

The left side of \eqref{fp:eq:four-state-pentagon} telescopes to
\(r_{32}r_{21}r_{10}G_0-G_3\).
The right side telescopes to \(r_{34}r_{40}G_0-G_3\).
The coefficient pentagon identifies these expressions before applying
them to any input.  Thus the boundary is exactly zero; a
degree-minus-two filling can be chosen to be zero for this choice
of edge homotopies.

The contractions \(h_T\) move a scalar from an overlap into
an endpoint module in which it is allowed.  They do not change it.
The extra primitive replaces \(\theta_c\) by \(e_1^c\)
only for coefficients already regular at \(c=0\) over \(S\).
All remaining operations evaluate components or insert units.
They introduce neither a new divisor nor a larger pole exponent.
The uniform output bounds therefore persist in the homotopies.
\end{proof}

The choice of \(G_T\) is part of this theorem.  Arbitrarily
chosen homotopies for the five edges need not satisfy
\eqref{fp:eq:four-state-pentagon}.  For the isolated three-leaf
rotation, this construction gives the explicit homotopy
\(i_R(e_1^xED_{13})-h_RY_R\): the left contraction is
annihilated by its first-endpoint projection.  Thus the same
normalization agrees with that three-input formula.

\subsubsection{The local ternary homotopies and a pentagon filling}

The local three-input homotopies admit a direct comparison with
\eqref{fp:eq:four-state-potentials}.  In the following table,
\(v\) is the retained local interval, \(w\) the inserted interval,
\(s\) the spectator interval, and \(X,L\) are odd state operators:
\[
\begin{array}{c|ccc|c|c}
 T\to U&v&w&s&X&L\\\hline
 0\to1&b&d&c&D_{13}&D_{14}+D_{24}+D_{34}\\
 1\to2&c&f&b&D_{14}&D_{23}\\
 2\to3&c&e&f&D_{24}&D_{12}+D_{13}+D_{14}\\
 0\to4&c&e&a&D_{14}+D_{24}&D_{12}\\
 4\to3&e&f&c&D_{13}+D_{14}&D_{34}
\end{array}
\]
Let \(h_w\) contract only the inserted interval, with its tensor
sign, and define
\[
 H^{\mathrm{loc}}_{UT}=e_1^vEX(1+\theta_sL)-h_wY_U.
\]
The placement of coefficients and operators follows the same graded
rule as \eqref{fp:eq:four-state-formulas}.  Deleting the old interval
from \(Y_T\) and the new interval from \(Y_U\) leaves the
difference \(\theta_vEX(1+\theta_sL)\).  The table verifies
this directly in each of its five rows.  Every coefficient of \(X\)
is in \(S\), so the displayed primitive is legal.  Hence
\[
 \partial H^{\mathrm{loc}}_{UT}=r_{UT}Y_T-Y_U.
\]
These are the local ternary homotopies with their unchanged spectator.
Their comparison with the coherent choices is explicit.  Set
\[
 Z_{UT}=H^{\mathrm{loc}}_{UT}-H_{UT},\qquad
 B_{UT}=h_UZ_{UT}.
\]
The map \(Z_{UT}\) is closed.  Full first-endpoint projection
of \(H^{\mathrm{loc}}_{UT}\) is, in the table's order,
\[
 0,\quad e_1^cED_{14},\quad e_1^cED_{24},\quad
 e_1^cE(D_{14}+D_{24}),\quad0.
\]
These values equal \(p_UH_{UT}\), by the formulas for \(C_T\).
Thus \(p_UZ_{UT}=0\), and the contraction gives
\(\partial B_{UT}=Z_{UT}\).
The degree-minus-two map
\[
 \mathcal K=r_{32}r_{21}B_{10}+r_{32}B_{21}+B_{32}
                -r_{34}B_{40}-B_{34}
\]
therefore satisfies the pentagon boundary identity
\begin{equation}\label{fp:eq:four-local-pentagon-filling}
 \partial\mathcal K=
 r_{32}r_{21}H^{\mathrm{loc}}_{10}+r_{32}H^{\mathrm{loc}}_{21}
       +H^{\mathrm{loc}}_{32}-r_{34}H^{\mathrm{loc}}_{40}-H^{\mathrm{loc}}_{34}.
\end{equation}
Indeed, substitute \(\partial B=H^{\mathrm{loc}}-H\) and use
\eqref{fp:eq:four-state-pentagon} to cancel the coherent terms.
This filling preserves the same pole tuples.  It proves coherence
for the specified local ternary homotopies, without assuming that
their two path sums are equal as degree-minus-one maps.

\begin{example}[A nonzero local pentagon boundary]
Take \(k=\mathbb Q\), \(a=2,b=3,c=5\), and inputs
\((x_0,x_0,\eta_0,\eta_0)\).  Thus \(d=5,e=8,f=10\).
The oriented local boundary on the right of
\eqref{fp:eq:four-local-pentagon-filling} is
\[
 -\frac{\theta_ce_1^f+e_1^c\theta_f}{40}
 -\frac{e_1^e\theta_f+\theta_ee_1^f}{120}.
\]
Its four summands are distinct interval-basis tensors, so the
boundary is nonzero.  The filling on these inputs is
\[
 \mathcal K=-\frac{e_1^ce_1^f}{40}-\frac{e_1^ee_1^f}{120}.
\]
Differentiating gives the displayed boundary.  Both contractions
consume all four generators, so the output state is \(1\).
This example distinguishes equality of local edge sums from the
chain-homotopy boundary identity.
\end{example}

\subsection{Mode coefficients and finite bars}
\label{fp:sec:four-mode-bars}

The coefficient diagram and its pentagon also act on associative
operator algebras.  This statement does not require their normal
products to be associative.  Define the finite mode algebra
\(\mathscr W_\kappa\) over \(k\) by
\[
 [p_m,q_n]=\kappa m\delta_{mn}1,\qquad
 [p_m,p_n]=[q_m,q_n]=0,\qquad m,n\ge1.
\]
Elements are finite linear combinations of finite words.  On
\(\mathsf F=k[h_0,h_1,\ldots]\), let
\(q_m=h_{m-1}\cdot\) and
\(p_m=\kappa m\partial_{h_{m-1}}\).
These actions satisfy the relations in every characteristic.
Ordered monomials \(q^\alpha p^\beta\) are linearly independent:
represent them on \(k[Q_m,P_m:m\ge1]\) with
\(q_m=Q_m\cdot\) and
\(p_m=\kappa m\partial_{Q_m}+P_m\cdot\).
Their values on \(1\) are the distinct monomials
\(Q^\alpha P^\beta\).  Commuting adjacent generators proves
spanning.  Thus a mode remains an algebra element even if its
particular Fock action vanishes in positive characteristic.

Give \(h_j\) weight \(j+1\) and let
\(\mathsf M=\bigoplus_w\mathsf F_w^\vee\).
The right action is \((\lambda A)(v)=\lambda(Av)\).
Every finite mode word has a fixed weight shift, so this action
preserves \(\mathsf M\).  Set
\(D_R=R\otimes_k\mathscr W_\kappa\) and extend these
module ends to each endpoint ring.  Each coefficient algebra acts
on the ends through first-endpoint evaluation of every interval.
No character of \(\mathscr W_\kappa\) is used.
For \(\kappa\ne0\), the relation
\([p_1,q_1]=\kappa1\) forbids such a character.

There is also an actual endomorphism carrier.  Let
\[
 \mathscr E=\bigoplus_{r\in\mathbb Z}
       \prod_{w\ge0}\operatorname{Hom}_k(\mathsf F_w,\mathsf F_{w+r}),
 \qquad\mathsf F_v=0\ (v<0).
\]
An element has finitely many weight shifts.  Composition makes
\(\mathscr E\) an associative algebra, since finite sets of
shifts add.  It acts on \(\mathsf F\) and preserves the
restricted dual \(\mathsf M\).  The mode representation maps
\(\mathscr W_\kappa\to\mathscr E\).  The constructions below
are natural under this map and keep each action unchanged.
If the full algebra \(\operatorname{End}_k(\mathsf F)\) is
used instead, its natural right end is the full algebraic dual
\(\mathsf F^*\); arbitrary endomorphisms need not preserve
\(\mathsf M\).

For an augmented dg algebra \(E\) over a current endpoint ring
\(L\), put \(\bar E=\ker\varepsilon_E\) and
\[
 B_L(E)=\bigoplus_{n\ge0}(s\bar E)^{\otimes_L n},
 \qquad |sv|=|v|-1.
\]
The bar differential has local components
\(b_1[v]=[\partial v]\) and
\(b_2[v|w]=(-1)^{|v|+1}[vw]\).
An operation starting at position \(j\) carries the prefix sign
\((-1)^{\sum_{i<j}(|v_i|-1)}\).
For genuine degree-zero module ends \(M,N\), use instead
\[
 B_L(M,E,N)=\bigoplus_{n\ge0}
 M\otimes_L(sE)^{\otimes_L n}\otimes_LN.
\]
The left outer face has positive sign.  The right outer face has
sign \((-1)^{\sum_i(|v_i|-1)}\).  The ends act through the
specified augmentations of the coefficient factors and the actual
state actions.  The state algebra itself needs no augmentation.
Associativity, the product rule, and the module identities cancel
the pairs of faces in \(b^2\).

\begin{theorem}[The pentagon with operator coefficients and module ends]
\label{fp:thm:four-bar-pentagon}
The coefficient pentagon induces a strict pentagon on reduced bars
when the state algebra is augmented.  It also induces a strict
pentagon on the two-sided bars with the stated mode or endomorphism
ends.  All maps act identically on state factors and module ends.
Over \(S\), each edge is a chain homotopy equivalence on these
finite-word bars.  The maps preserve cuts, word length, and total
spatial pole bounds.  Reduced bars use deconcatenation.  Two-sided
bars use balanced cuts through the regular state module.
\end{theorem}

\begin{proof}
A strict algebra chain map acts entry by entry and commutes with
all internal and multiplication faces.  Its preservation of the
endpoint action gives the same assertion for outer faces.
Consequently \eqref{fp:eq:four-pentagon-coefficients} gives the strict
bar pentagon in every length, with the identity on both ends.

For one deleted invertible interval, its tensor homotopy \(h\)
and projection \(F\) satisfy
\[
 h(vw)=h(v)w+(-1)^{|v|}F(v)h(w).
\]
For one interval this follows from the four endpoint actions.
Tensoring with the other factors gives the displayed Koszul sign.
A homotopy on words is
\[
 K[v_1|\cdots|v_n]=
 \sum_{j=1}^n(-1)^{\sum_{i<j}(|v_i|-1)}
 [Fv_1|\cdots|Fv_{j-1}|hv_j|v_{j+1}|\cdots|v_n].
\]
The internal terms of \(bK+Kb\) telescope to \(1-B(F)\).
The adjacent multiplication terms reduce to
\[
 (-1)^{|v|}[h(v)w]+[F(v)h(w)]
       +(-1)^{|v|+1}[h(vw)]=0.
\]
The outer terms cancel because \(h\) has zero first endpoint
and \(F\) preserves endpoint evaluation.  This proves the
homotopy-equivalence assertion for each edge and its reverse.

A cut before or after an entrywise map gives the same two
subwords.  For two-sided bars, insert the regular state module \(D\)
at a cut and balance over \(D\):
\[
 \lambda[v_1|\cdots|v_n]w\longmapsto
 \sum_j\lambda[v_1|\cdots|v_j]1
 \otimes_D1[v_{j+1}|\cdots|v_n]w.
\]
The two adjacent middle faces cancel after balancing, with opposite
prefix signs.  All other faces belong to one side of the cut.
Endpoint rotations preserve these actions and the inserted unit.
Hence they commute with balanced cuts.  Two successive cuts choose
three consecutive subwords, independently of their order.

For pole filtrations, use the scalar lattices
\(\prod_\delta\delta^{-N_\delta}R\), intersected with each
coefficient module, and a specified regular state lattice.  On words
take the linear span of tensors whose tuples sum to at most the
stated bound.  Multiplication adds tuples.  The maps and the displayed
homotopy preserve every scalar coefficient, proving the assertion.
\end{proof}

This proof contracts one changed interval at each edge.  The tensor
homotopy \eqref{fp:eq:four-tensor-contraction} used for state outputs
is only a chain homotopy; no multiplicative property of that sum
is required.  For a projection deleting several intervals on bars,
one may compose the single-interval contractions just proved.
All words remain finite.  No evaluation of an infinite distribution
of mode words is part of the construction.

\subsubsection{The grouping comparison on every finite word}

The nonmultiplicative map \(\mu_T\) also has explicit higher
bar components.  The following construction applies to a fixed tree
over \(L=R_{J_T}\).  Write \(C=\mathcal C_{J_T}(D_R)\)
and \(E=E_T(D_R)\otimes_RL\).  Use the contraction
\((j,\mu,H)\) from \eqref{fp:eq:four-grouping-contraction}.
On tensor words set \(I_0=B(j)\), let \(P_0\) apply
\(\mu\) to each entry, and put \(F_0=I_0P_0\).
For the internal differentials alone, the tensor homotopy is
\[
 H_0[v_1|\cdots|v_n]=
 \sum_{r=1}^n(-1)^{\sum_{i<r}(|v_i|-1)}
 [j\mu v_1|\cdots|j\mu v_{r-1}|Hv_r|v_{r+1}|\cdots|v_n].
\]
It satisfies
\[
 b_1H_0+H_0b_1=1-F_0,\qquad
 H_0^2=H_0I_0=P_0H_0=0.
\]
Let \(\delta\) be the multiplication part of the source bar
differential.  For two-sided bars include the two outer faces in
\(\delta\).  In either case \(\delta\) lowers length.
Define, on a word of length \(n\),
\begin{equation}\label{fp:eq:four-grouping-bar}
 P=P_0\sum_{r=0}^{n}(-\delta H_0)^r,
 \qquad
 \mathcal H=H_0\sum_{r=0}^{n}(-\delta H_0)^r.
\end{equation}
These are finite formulas with no convergence requirement.

\begin{theorem}[Higher corrections for grouping intervals]
\label{fp:thm:four-grouping-bar}
The map \(P\) in \eqref{fp:eq:four-grouping-bar} is a chain
homotopy inverse to the strict bar inclusion \(I_0\).
Its length-one component is \(\mu_T\), and its target has the
ordinary algebra bar differential, with the stated module ends
when present.  On reduced bars, \(P\) is a coaugmented dg
coalgebra map.  The two-sided maps commute with balanced cuts.
The construction preserves maximal word length and total pole tuples.
It applies to the actual mode and endomorphism coefficients above.
\end{theorem}

\begin{proof}
Write \(d=b_1\), \(Q=1+\delta H_0\), and let \(\delta_E\)
be the multiplication and applicable outer faces of the target.
The strict algebra inclusion preserves endpoint actions, so
\(\delta I_0=I_0\delta_E\).  Consequently
\(H_0\delta I_0=0\), and no higher target multiplication is needed.
The inverse \(Q^{-1}\) is the finite series in
\eqref{fp:eq:four-grouping-bar}.  The equations
\(d\delta+\delta d=0\), \(\delta^2=0\), and
\(dH_0+H_0d=1-F_0\) give
\[
 dQ-Qd=-\delta(1-F_0),\qquad
 dQ^{-1}-Q^{-1}d=Q^{-1}\delta(1-F_0)Q^{-1}.
\]
It follows that
\[
 \begin{aligned}
 (d+\delta)\mathcal H+\mathcal H(d+\delta)
 &=1-I_0P+H_0Q^{-1}\delta(\delta H_0+F_0)Q^{-1}\\
 &=1-I_0P.
 \end{aligned}
\]
The last term is zero: \(\delta^2=0\) and
\(\delta F_0=I_0\delta_EP_0\), while
\(H_0Q^{-1}I_0=H_0I_0=0\).
Also \(PI_0=1\).  Commuting the last homotopy identity with
the source differential shows that \(I_0P\) is a chain map.
Since \(I_0\) is a split injection and a chain map to the
ordinary target bar, \(P\) is a chain map to that bar.
This proves the homotopy-equivalence assertion.
The same identities hold with the two outer faces: their square
and their commutation with the internal differential follow from
the module identities, and \(I_0\) preserves their actions.

Here is a direct check of the coalgebra assertion.  The tensor formula
has
\[
 \Delta H_0=(F_0\otimes H_0+H_0\otimes1)\Delta,
 \qquad \Delta\delta=(\delta\otimes1+1\otimes\delta)\Delta.
\]
Moreover \(PH_0=0\), \(PF_0=P_0\), and
\(P=P_0-P\delta H_0\).  Induct on input length.
The induction applies to \(\delta H_0\), since it lowers
length.  The two terms containing \(PH_0\) vanish, leaving
\[
 \begin{aligned}
 \Delta P
 &=\bigl(P_0\otimes P_0-P_0\otimes P\delta H_0
                  -P\delta H_0\otimes P\bigr)\Delta\\
 &=(P\otimes P)\Delta.
 \end{aligned}
\]
The empty word gives the base case.  For module ends, use this
same calculation for the family of bars \(B(M,C,N)\), with
\(B(M,C,D)\otimes_D B(D,C,N)\) at a cut.  The middle-face
cancellation proved above gives the displayed differential identity
for that balanced cut.  The tensor-homotopy identity separates terms
according to the side on which the homotopy acts, with its original prefix sign.
All maps are \(D\)-linear at the inserted ends and preserve
their units, so the calculation descends to the balanced tensor.

The linear component is \(\mu\): internal multiplication
cannot act on one reduced entry.  For a length-one two-sided
bar, each outer face of \(H_0[v]\) is zero because \(H\)
has zero first-endpoint value.  Thus the same statement holds there.
Finally every factor in the finite series changes only interval
coordinates or multiplies existing spatial coefficients in order.
Word length decreases under \(\delta\), and multiplication
adds the original pole bounds.  This proves the filtration assertions.
\end{proof}

The strict reorderings \(\rho\) give the full-divisor bar pentagon.
The maps \(P\) compare its vertices with the native mode and
coefficient bars on their exact localization domains.  On the linear
state component they give \(\mu_T\mathcal U=Y_T\).
Thus the common state map and the native bar multiplication have an
explicit comparison, including the higher terms forced by grouping.

\subsubsection{Compatibility of the comparison maps around the pentagon}

Over \(S\), identify all enriched source bars with the common
bar \(\mathcal B\) of \(\mathcal C(D_R)\).
Let \(\Pi:\mathcal B\to B(A_c\otimes D_R)\) evaluate
the five other intervals, and let \(I_c\) insert their units.
Use the corresponding two-sided notation when module ends are present.
The single-interval bar contractions above give an explicit
\(K_c\) with \(bK_c+K_cb=1-I_c\Pi\).
For example, if their order is \(a,b,d,e,f\), take
\[
 K_c=K_a+B(F_a)K_b+B(F_aF_b)K_d
       +B(F_aF_bF_d)K_e+B(F_aF_bF_dF_e)K_f.
\]
The chain equations telescope.  Each \(P_T\) from
\eqref{fp:eq:four-grouping-bar} satisfies
\(P_TI_c=B(i_T)\): indeed, \(I_c=B(j_T)B(i_T)\)
and \(P_TB(j_T)=1\).
Define comparison homotopies
\[
 L_T=P_TK_c,\qquad
 A_{UT}=B(r_{UT})L_T-L_U.
\]
Their chain equations are
\[
 bL_T+L_Tb=P_T-B(i_T)\Pi,
 \qquad bA_{UT}+A_{UT}b=B(r_{UT})P_T-P_U.
\]
The terms through \(\Pi\) cancel because
\(r_{UT}i_T=i_U\).  Transporting the \(A_{UT}\) around
the two pentagon paths gives respectively
\(B(r_{32}r_{21}r_{10})L_0-L_3\) and
\(B(r_{34}r_{40})L_0-L_3\).  They are identical by
\eqref{fp:eq:four-pentagon-coefficients}.
Thus the native bars, the common enriched bar, and their explicit
comparison maps have compatible pentagon homotopies on the common open.
All homotopies are finite on every word and preserve pole tuples.

\subsection{Boundary and completion limits}
\label{fp:sec:four-boundary-limits}

There is a common formal domain retaining the same \(c\)-collision.
Introduce variables \(A,B,\lambda,\epsilon\) and put
\[
 U=k[A,B,(AB(A+B))^{-1}][[\epsilon,\lambda]][\lambda^{-1}],
\]
\[
 a=A\lambda,\qquad b=B\lambda,\qquad c=\epsilon\lambda.
\]
Then \(d=(A+B)\lambda\), \(e=(B+\epsilon)\lambda\),
and \(f=(A+B+\epsilon)\lambda\).
The constant terms \(B\) and \(A+B\) are units in the
coefficient ring, so these formulas define
\(S\to U\) and \(S[c^{-1}]\to U[\epsilon^{-1}]\).
The polynomial substitution is injective.  Localization preserves
this injectivity because all named denominators become units.
A pole tuple \(N\) maps into
\[
 \epsilon^{-N_c}\lambda^{-\sum_{\delta\in\mathcal D}N_\delta}
 k[A,B,(AB(A+B))^{-1}][[\epsilon,\lambda]].
\]
The numerator-one monomial attains both lower exponents.
The weightwise state formulas and their homotopies therefore extend
to this specified domain for polynomial inputs.

An arbitrary Laurent series in a tree target need not admit this
substitution.  For example,
\(\sum_{n\ge0}a^nc^{-n^2}\) belongs to
\(k((c))((b))((a))\), but its substituted terms have
unbounded negative \(\epsilon\)- and \(\lambda\)-orders.
There is no element of \(U[\epsilon^{-1}]\) with that expansion.
Likewise, in the product of state weights, \(\eta_j\to0\),
but the scalar coefficient of \(Y_c(x_0,\eta_j)\) is
\(\theta_c/c^{j+1}\).  It does not tend to zero in a discrete
scalar-state coefficient space.  Thus completed inputs require
additional topology and continuity statements.

The all-divisor diagram \eqref{fp:eq:four-common-carrier} retains
the original boundary data.  Its localized tree equivalences and
its state homotopies have the precise domains proved above.
The coefficient pentagon holds on the original ring as an identity
of maps; its edges are not equivalences there.
A common geometric factorization object with equivalences across
all six collision divisors would require different fibers and a
proved comparison to this diagram.  The nonzero local class in
Theorem~\ref{fp:thm:four-coefficient-pentagon} rules out obtaining
such an object by identifying the five unchanged scalar trees.
